% GENERATED FILE. Edit canonical lessons, then run npm run generate:books.
% canonical-source-hash: fnv1a64:3e01857ff9dd3dac
% canonical-lessons: TE-C364-banisa, TE-C364-vimukti, TE-C364-poratam, TE-C364-tyagam, TE-C364-virudu

\chapter{The Price of Freedom}
\label{ch:te-the-price-of-freedom}
\hlchaptermodality{364}

\begin{chapteropening}
\textbf{By the end of this chapter:} \emph{I can say a slave, liberation, a struggle, a sacrifice and a brave man.}

Complete the last lesson of chapter 364: I can say a slave, liberation, a struggle, a sacrifice and a brave man.
\end{chapteropening}

\section[\texorpdfstring{\te{బానిస} (bānisa)}{బానిస (bānisa)}]{\te{బానిస} (bānisa) — a slave}

\label{lesson:TE-C364-banisa}



\begin{quote}
Before the new one: say the Telugu for to incite, then the Telugu for independence.
\end{quote}

\subsection*{You'll want to know: \te{బానిస}}
\textbf{\te{బానిస}} — \emph{bānisa} — \textquotedblleft{}a slave\textquotedblright{}.

\begin{quote}
\te{పూర్వం} \te{బానిసలను} \te{అమ్మేవారు}.
\end{quote}

\begin{quote}
\emph{pūrvaṁ bānisalanu ammēvāru.}
\end{quote}

\begin{quote}
In the past, slaves were sold.
\end{quote}

Plural: \te{బానిసలు}, \emph{bānisalu}.

\te{బానిసత్వం} is slavery.

\begin{cousinweb}
\textbf{\te{బానిస}} (\emph{bānisa}) is the Telugu word for a slave; its origin is not settled.
\end{cousinweb}

\begin{grammarlens}[title={What you've built}]
One more word for telling what happened.
\end{grammarlens}

\subsection*{Your turn}
\begin{itemize}
\raggedright
  \item \emph{Say it:} \emph{bānisa}
  \item \emph{Say it:} \emph{bānisa}, once more
  \item \emph{Say it:} \emph{pūrvaṁ bānisalanu ammēvāru.} — In the past, slaves were sold.
  \item \emph{Recall:} say the Telugu for to incite, then the Telugu for independence, then say \emph{bānisa} again
\end{itemize}

\begin{tcolorbox}[breakable,colback=teal!4,colframe=teal!35!black,title={Before you move on}]
What does \te{బానిస} mean? (\textquotedblleft{}A slave\textquotedblright{}.) Say it once more.
\end{tcolorbox}
\section[\texorpdfstring{\te{విముక్తి} (vimukti)}{విముక్తి (vimukti)}]{\te{విముక్తి} (vimukti) — liberation; release}

\label{lesson:TE-C364-vimukti}



\begin{quote}
Before the new one: say the Telugu for independence, then the Telugu for a slave.
\end{quote}

\subsection*{You'll want to know: \te{విముక్తి}}
\textbf{\te{విముక్తి}} — \emph{vimukti} — \textquotedblleft{}liberation; release\textquotedblright{}.

\begin{quote}
\te{బానిసలకు} \te{చివరికి} \te{విముక్తి} \te{లభించింది}.
\end{quote}

\begin{quote}
\emph{bānisalaku civariki vimukti labhiñcindi.}
\end{quote}

\begin{quote}
At last the slaves won their freedom.
\end{quote}

\begin{cousinweb}
\textbf{\te{విముక్తి}} (\emph{vimukti}) is borrowed from Sanskrit \emph{vimukti}, \textquotedblleft{}giving up; release\textquotedblright{} (Monier-Williams), from \emph{vi-muc}, \textquotedblleft{}to set loose\textquotedblright{}.
\end{cousinweb}

\begin{grammarlens}[title={What you've built}]
One more word for telling what happened.
\end{grammarlens}

\subsection*{Your turn}
\begin{itemize}
\raggedright
  \item \emph{Say it:} \emph{vimukti}
  \item \emph{Say it:} \emph{vimukti}, once more
  \item \emph{Say it:} \emph{bānisalaku civariki vimukti labhiñcindi.} — At last the slaves won their freedom.
  \item \emph{Recall:} say the Telugu for independence, then the Telugu for a slave, then say \emph{vimukti} again
\end{itemize}

\begin{tcolorbox}[breakable,colback=teal!4,colframe=teal!35!black,title={Before you move on}]
What does \te{విముక్తి} mean? (\textquotedblleft{}Liberation; release\textquotedblright{}.) Say it once more.
\end{tcolorbox}
\section[\texorpdfstring{\te{పోరాటం} (pōrāṭaṁ)}{పోరాటం (pōrāṭaṁ)}]{\te{పోరాటం} (pōrāṭaṁ) — a struggle}

\label{lesson:TE-C364-poratam}



\begin{quote}
Before the new one: say the Telugu for a slave, then the Telugu for liberation.
\end{quote}

\subsection*{You'll want to know: \te{పోరాటం}}
\textbf{\te{పోరాటం}} — \emph{pōrāṭaṁ} — \textquotedblleft{}a struggle\textquotedblright{}.

\begin{quote}
\te{స్వాతంత్ర్యం} \te{కోసం} \te{పోరాటం} \te{చాలా} \te{కాలం} \te{జరిగింది}.
\end{quote}

\begin{quote}
\emph{svātantryaṁ kōsaṁ pōrāṭaṁ cālā kālaṁ jarigindi.}
\end{quote}

\begin{quote}
The struggle for independence went on for a long time.
\end{quote}

Plural: \te{పోరాటాలు}, \emph{pōrāṭālu}.

The noun of \te{పోరాడు}, to fight.

\begin{cousinweb}
\textbf{\te{పోరాటం}} (\emph{pōrāṭaṁ}) is native Dravidian, the noun of \emph{pōru}, \textquotedblleft{}to fight, struggle\textquotedblright{} (DEDR 4540).
\end{cousinweb}

\begin{grammarlens}[title={What you've built}]
One more word for telling what happened.
\end{grammarlens}

\subsection*{Your turn}
\begin{itemize}
\raggedright
  \item \emph{Say it:} \emph{pōrāṭaṁ}
  \item \emph{Say it:} \emph{pōrāṭaṁ}, once more
  \item \emph{Say it:} \emph{svātantryaṁ kōsaṁ pōrāṭaṁ cālā kālaṁ jarigindi.} — The struggle for independence went on for a long time.
  \item \emph{Recall:} say the Telugu for a slave, then the Telugu for liberation, then say \emph{pōrāṭaṁ} again
\end{itemize}

\begin{tcolorbox}[breakable,colback=teal!4,colframe=teal!35!black,title={Before you move on}]
What does \te{పోరాటం} mean? (\textquotedblleft{}A struggle\textquotedblright{}.) Say it once more.
\end{tcolorbox}
\section[\texorpdfstring{\te{త్యాగం} (tyāgaṁ)}{త్యాగం (tyāgaṁ)}]{\te{త్యాగం} (tyāgaṁ) — a sacrifice}

\label{lesson:TE-C364-tyagam}



\begin{quote}
Before the new one: say the Telugu for liberation, then the Telugu for a struggle.
\end{quote}

\subsection*{You'll want to know: \te{త్యాగం}}
\textbf{\te{త్యాగం}} — \emph{tyāgaṁ} — \textquotedblleft{}a sacrifice\textquotedblright{}.

\begin{quote}
\te{అమ్మ} \te{పిల్లల} \te{కోసం} \te{చాలా} \te{త్యాగాలు} \te{చేసింది}.
\end{quote}

\begin{quote}
\emph{amma pillala kōsaṁ cālā tyāgālu cēsindi.}
\end{quote}

\begin{quote}
Mother made many sacrifices for her children.
\end{quote}

\te{త్యాగం} \te{చేయు}, to make a sacrifice.

\begin{cousinweb}
\textbf{\te{త్యాగం}} (\emph{tyāgaṁ}) is borrowed from Sanskrit \emph{tyāga}, \textquotedblleft{}abandonment\textquotedblright{}, and later \textquotedblleft{}liberality\textquotedblright{} (CDIAL 5985): giving something up.
\end{cousinweb}

\begin{grammarlens}[title={What you've built}]
One more word for telling what happened.
\end{grammarlens}

\subsection*{Your turn}
\begin{itemize}
\raggedright
  \item \emph{Say it:} \emph{tyāgaṁ}
  \item \emph{Say it:} \emph{tyāgaṁ}, once more
  \item \emph{Say it:} \emph{amma pillala kōsaṁ cālā tyāgālu cēsindi.} — Mother made many sacrifices for her children.
  \item \emph{Recall:} say the Telugu for liberation, then the Telugu for a struggle, then say \emph{tyāgaṁ} again
\end{itemize}

\begin{tcolorbox}[breakable,colback=teal!4,colframe=teal!35!black,title={Before you move on}]
What does \te{త్యాగం} mean? (\textquotedblleft{}A sacrifice\textquotedblright{}.) Say it once more.
\end{tcolorbox}
\section[\texorpdfstring{\te{వీరుడు} (vīruḍu)}{వీరుడు (vīruḍu)}]{\te{వీరుడు} (vīruḍu) — a brave man; a hero}

\label{lesson:TE-C364-virudu}



\begin{quote}
Before the new one: say the Telugu for a struggle, then the Telugu for a sacrifice.
\end{quote}

\subsection*{You'll want to know: \te{వీరుడు}}
\textbf{\te{వీరుడు}} — \emph{vīruḍu} — \textquotedblleft{}a brave man; a hero\textquotedblright{}.

\begin{quote}
\te{ఆ} \te{వీరుడు} \te{ఒంటరిగా} \te{శత్రువులతో} \te{పోరాడాడు}.
\end{quote}

\begin{quote}
\emph{ā vīruḍu oṇṭarigā śatruvulatō pōrāḍāḍu.}
\end{quote}

\begin{quote}
That brave man fought the enemies alone.
\end{quote}

Plural: \te{వీరులు}, \emph{vīrulu}.

\begin{cousinweb}
\textbf{\te{వీరుడు}} (\emph{vīruḍu}) is borrowed from Sanskrit \emph{vīra}, \textquotedblleft{}a man, a hero\textquotedblright{} (CDIAL 12056), with the Telugu masculine ending \emph{-uḍu}.
\end{cousinweb}

\begin{grammarlens}[title={What you've built}]
One more word for telling what happened.
\end{grammarlens}

\subsection*{Your turn}
\begin{itemize}
\raggedright
  \item \emph{Say it:} \emph{vīruḍu}
  \item \emph{Say it:} \emph{vīruḍu}, once more
  \item \emph{Say it:} \emph{ā vīruḍu oṇṭarigā śatruvulatō pōrāḍāḍu.} — That brave man fought the enemies alone.
  \item \emph{Recall:} say the Telugu for a struggle, then the Telugu for a sacrifice, then say \emph{vīruḍu} again
\end{itemize}

\begin{tcolorbox}[breakable,colback=teal!4,colframe=teal!35!black,title={Before you move on}]
What does \te{వీరుడు} mean? (\textquotedblleft{}A brave man; a hero\textquotedblright{}.) Say it once more.
\end{tcolorbox}
